% TOP_PROBLEM: {"id": 2609, "title": "Kourovka Notebook Problem 21.100", "category_id": 17, "category": {"id": 17, "name": "group_theory", "display_name": "Group Theory", "description": "Problems about groups, group actions, representations, and related algebraic structures.", "slug": "group-theory", "order_index": 17, "created_at": "2026-07-31T15:26:25.671Z"}, "status": "open", "problem_number": "KOU-21.100", "set_id": 10}
% FIRST_POSTED: 2026-10-01
% ENTRY_KIND: statement_only
% UnsolvedMath ID 2609; source catalogue code KOU-21.100
% !TeX program = lualatex
% Definition and sources only; this file is not an LLM solution attempt.
\documentclass[11pt,a4paper]{article}
\usepackage[margin=25mm]{geometry}
\usepackage{fontspec}
\setmainfont{lmroman10-regular.otf}[BoldFont=lmroman10-bold.otf,
 ItalicFont=lmroman10-italic.otf,BoldItalicFont=lmroman10-bolditalic.otf]
\usepackage{amsmath,amssymb,mathtools}
\usepackage{unicode-math}
\setmathfont{latinmodern-math.otf}
\AtBeginDocument{\let\setminus\smallsetminus}
\usepackage{xurl,hyperref}
\hypersetup{colorlinks=true,urlcolor=blue}
\setcounter{secnumdepth}{0}
\setlength{\parindent}{0pt}
\setlength{\parskip}{5pt}
\setlength{\emergencystretch}{3em}
\begin{document}
\section*{Kourovka Notebook Problem 21.100}\label{2609}
{\small UnsolvedMath ID 2609 · KOU-21.100 · Group Theory · Kourovka Notebook - New Problems, Issue 21}
\subsection{Definitions and mathematical statement}
Let $A$ and $G$ be finite groups, with an action $A\to\operatorname{Aut}(G)$ satisfying $\gcd(|A|,|G|)=1$. This is a coprime action; it need not be faithful. Write $a(g)$ for the image of $g$ under $a$, and define the common fixed-point subgroup
\[
 C=C_G(A)=\{g\in G:a(g)=g\text{ for every }a\in A\}.
\]
Let $C'=[C,C]$ be the subgroup generated by the commutators of elements of $C$.

An irreducible complex character of $G$ is the trace function of a finite-dimensional irreducible complex representation, taken once for each representation isomorphism class. Let $\operatorname{Irr}(G)$ be the set of these characters. A character $\chi$ is $A$-invariant if $\chi(a(g))=\chi(g)$ for every $a\in A$ and $g\in G$. Its restriction to $C$ is nowhere zero if $\chi(c)\ne0$ for every $c\in C$.

Does every such coprime action satisfy the exact counting identity
\[
 \left|\{\chi\in\operatorname{Irr}(G):
 \chi\text{ is }A\text{-invariant and }\chi(c)\ne0\ (c\in C)\}\right|
 =|C/C'|?
\]
The restriction need not itself be irreducible. The count includes the principal character and counts distinct irreducible characters even when their degrees agree. The right side is the order of the abelianization of the fixed-point subgroup.
\subsection{Short English statement}
Under coprime automorphisms, does the number of invariant irreducible characters nonzero everywhere on the fixed subgroup equal that subgroup’s abelianization order?
\subsection{Sources}
\begin{itemize}
\item[S1] ULAM.AI and the UnsolvedMath Contributors, supplied problems.json, record 2609 (KOU-21.100), Kourovka Notebook - New Problems, Issue 21. Catalogue identity and source transcription; CC BY 4.0.
\url{https://huggingface.co/datasets/ulamai/UnsolvedMath}.
\item[S2] The Kourovka Notebook, 21st edition, version 47 (30 September 2026), new problems, Problem 21.100. Expanded formulation of the supplied catalogue statement.
\url{https://arxiv.org/pdf/1401.0300v47}.
\end{itemize}
\end{document}
